% !TITLE 诗经·秦风·无衣
% !AUTHOR 无名氏
% !DYNASTY 先秦
% !GENRE 四言诗
% !ORDER 8

\begin{poem}
岂曰无衣？与子同袍。
王于兴师，修我戈矛。与子同仇！

岂曰无衣？与子同泽。
王于兴师，修我矛戟。与子偕作！

岂曰无衣？与子同裳。
王于兴师，修我甲兵。与子偕行！
\end{poem}

\begin{appreciation}
《无衣》是《诗经》中最慷慨激昂的战歌，也是中国文学史上最早的同仇敌忾之声。

诗的背景，传统认为是秦人抗击西戎的战争。秦国地处西陲，长期与戎狄交战，民风剽悍，尚武精神浓厚。这首诗不是某一个人的独白，而是一群战士的合唱——那种一呼百应的气势，从字里行间扑面而来。

"岂曰无衣？与子同袍"——谁说没有衣服？我和你同穿一件战袍。这不是在讨论穿衣问题，而是在表达一种生死与共的情谊。在战场上，袍泽之情是最珍贵的。"同袍"这个词，后来成为战友的代称，就是从这里来的。

三章的层次逐渐递进。第一章"与子同仇"——我们共同面对仇敌，这是态度的宣示；第二章"与子偕作"——我们一起行动起来，这是行动的号召；第三章"与子偕行"——我们一同奔赴战场，这是决绝的誓言。从"同仇"到"偕作"再到"偕行"，情感越来越激昂，节奏越来越快，像战鼓一声紧似一声。

这首诗的语气是粗粝的、直接的，没有修饰，没有婉转。但正是这种不加修饰的真诚，让它具有了震撼人心的力量。它不是一个人在抒情，而是一个民族在呐喊。读这首诗，你能感受到三千年前秦地男儿的热血——那种面对外敌入侵时，不分你我、不计生死的集体意志。

后世引用这首诗，往往取其"同仇敌忾"之意。但诗中更动人的，是那种朴素的同伴情谊。"与子同袍"四个字，说的是战场上的相互扶持，也是人生路上最可贵的患难与共。
\end{appreciation}
